\item \textbf{Desplazamiento del máximo central por una lámina delgada.}
Su grupo está trabajando en un laboratorio de optoelectrónica. Una lámina de plástico transparente que tiene un índice de refracción $n$ y espesor $t$ se coloca entre la ranura superior y la pantalla de un experimento de Young, en una orientación tal que la luz pasa a través del plástico perpendicularmente a sus superficies (figura TP36.2). Cuando se hace esto, el máximo central del patrón de interferencia se mueve hacia arriba en la pantalla por una distancia $y'$. Investigue esta situación y determine una expresión para la distancia $y'$ en términos de la separación de rendijas $d$, la distancia a la pantalla $L$, el índice de refracción $n$ y el espesor $t$.